\begin{frame}{Updated Domain Model}
\begin{center}
\begin{tikzpicture}[node distance=.78cm,box/.style={draw,rounded corners,align=center,minimum width=2.15cm,minimum height=.7cm,fill=softgray},arr/.style={-{Stealth},thick}]
\node[box] (env) {Environment};
\node[box,fill=softred!18,right=of env] (ctx) {Missing/noisy\\context};
\node[box,right=of ctx] (pol) {\mab/\cmab/\icmab\\policy};
\node[box,right=of pol] (act) {Routing\\action};
\node[box,fill=softgreen!18,below=of act] (out) {Reward +\\fairness state};
\node[box,left=of out] (rep) {Reporting +\\validation};
\draw[arr] (env)--(ctx);\draw[arr] (ctx)--(pol);\draw[arr] (pol)--(act);\draw[arr] (act)--(out);\draw[arr] (out)--(rep);\draw[arr,dashed] (rep) to[bend right=35] (pol);
\end{tikzpicture}
\end{center}
\begin{itemize}
\item One policy sees only partial feedback, not every counterfactual.
\item Fairness risk appears when context is systematically worse for a group, site, flow, or route class.
\end{itemize}
\end{frame}

\begin{frame}{Architecture Overview}
\begin{center}
\begin{tikzpicture}[layer/.style={draw,rounded corners,align=center,minimum width=10.5cm,minimum height=.72cm,fill=softgray},arr/.style={-{Stealth},thick,deepblue}]
\node[layer,fill=softblue!18] (data) {Layer 1: Evidence, Data, and Context -- COVID evidence, clinical fixture, quantum states};
\node[layer,fill=softgreen!18,below=.3cm of data] (env) {Layer 2: Environments and Evaluators -- ClinicalEnvironment, quantum testbed, runners};
\node[layer,fill=ritorange!18,below=.3cm of env] (pol) {Layer 3: Policies and Allocation -- Baseline, FairAllocatorV1, ClinicalBanditUCB};
\node[layer,fill=softpurple!18,below=.3cm of pol] (rep) {Layer 4: EQUITAS, Reporting, and Validation -- fairness states, summaries, artifacts};
\draw[arr](data)--(env);\draw[arr](env)--(pol);\draw[arr](pol)--(rep);
\end{tikzpicture}
\end{center}
\begin{block}{Design principle}
Low coupling: each layer exposes records/artifacts to the next layer. High cohesion: each layer has one responsibility.
\end{block}
\end{frame}

\begin{frame}{Policy Comparison}
\begin{center}
\begin{tabular}{p{0.20\textwidth}p{0.31\textwidth}p{0.34\textwidth}}
\toprule
Policy family & Decision signal & Expected failure mode \\
\midrule
\mab & action history only & repeats historical allocation patterns \\
\cmab & observed context & improves utility but can inherit biased or incomplete context \\
\icmab & context plus missingness/fairness state & can learn when context itself is unreliable \\
\equitas-mediated & policy output plus disparity audit & constrains unsafe disparity while preserving utility \\
\bottomrule
\end{tabular}
\end{center}
\begin{equation*}
\text{decision}_t = \arg\max_a \left[\text{utility}_{t,a} - \lambda\,\text{expected disparity}_{t,a}\right]
\end{equation*}
\end{frame}

\begin{frame}{Preliminary Framework Results}
\small
\begin{center}
\begin{tabular}{p{0.42\textwidth}p{0.43\textwidth}}
\toprule
Evidence item & Current result \\
\midrule
Clinical embedded path & evaluator, runner, environment, allocator, models, and EQUITAS context run end-to-end \\
Clinical fairness smoke & 16 fairness states across four model families \\
Non-oracle efficiency & FairAllocatorV1: 91.23\%; ClinicalBanditUCB: 84.21\%; ClinicalBaseline: 75.09\% oracle efficiency \\
Reporting layer & 21 manifest-linked artifacts and model-stratified summaries \\
Quantum compatibility & legacy quantum path preserved with fairness sidecar artifacts \\
\bottomrule
\end{tabular}
\end{center}
\small Boundary: active fairness-adjusted quantum learning still requires an explicit policy-update mechanism and full validation.
\end{frame}

\begin{frame}{GitHub Code Review Break}
\begin{block}{Show 2--3 classes not used in the midterm}
Open GitHub and connect each class to the architecture diagram.
\end{block}
\small
\begin{center}
\begin{tabular}{p{0.32\textwidth}p{0.48\textwidth}}
\toprule
Class/file & Architecture connection \\
\midrule
\texttt{ClinicalExperimentEvaluator} & Layer 2 evaluator orchestration and artifact capture \\
\texttt{ClinicalExperimentRunner} & Layer 2 execution over seeds, models, and allocations \\
\texttt{FairAllocatorV1} or \texttt{ClinicalBanditUCB} & Layer 3 policy/allocation decision logic \\
\texttt{equitas\_mediator.py} & Layer 4 fairness audit and mitigation context \\
\bottomrule
\end{tabular}
\end{center}
\end{frame}

\begin{frame}{Demo Plan and Next Semester}
\begin{columns}[T]
\begin{column}{0.48\textwidth}
\textbf{Demo path}
\begin{enumerate}
\item run clinical fairness smoke/notebook
\item show fairness states and summaries
\item open report artifacts and manifests
\item point to COVID clinical foundation
\end{enumerate}
\end{column}
\begin{column}{0.48\textwidth}
\textbf{Next semester}
\begin{enumerate}
\item convert COVID cases into executable environments
\item complete active quantum fairness updates
\item expand validation hub
\item link results to paper-ready reporting
\end{enumerate}
\end{column}
\end{columns}
\end{frame}

\begin{frame}{Limitations and Takeaway}
\begin{itemize}
\item COVID analysis uses aggregate literature-grounded evidence, not patient-level data.
\item PCR/antigen false-positive disparity is not claimed; pulse oximetry is the defensible diagnostic case.
\item Fair CMAB allocations are counterfactual evaluation targets, not historical claims.
\end{itemize}
\begin{block}{Takeaway}
Fair sequential decision-making is not achieved by pretending context does not exist. It is achieved by measuring context, auditing disparities, and routing scarce resources where the cost of being wrong is highest.
\end{block}
\end{frame}
